\subsection{Stop Words and TF-IDF}
\label{sec:stop-words-tfidf}
\slidesrc{Syntax: RE \& word: 21--47}

\subsubsection{Stop words}
\slidesrc{Syntax: RE \& word: 21--23}

\begin{nbdef}{Stop words}
\textbf{Stop words} are words that are filtered out before or after processing natural
language text. The term usually refers to the \emph{most common} words of a language
(\emph{the}, \emph{of}, \emph{and}, \emph{a}, \emph{to}, \ldots), but there is no single
universal list used by all NLP tools, and some tools do not use one at all.
\end{nbdef}

\begin{itemize}[nosep]
  \item \textbf{How a stop list is built}: sort the terms by \emph{collection frequency}
        (the total number of times each term appears in the document collection), take
        the most frequent ones, and often filter them by hand according to their meaning
        in the domain of the documents. The words of the stop list are then discarded
        during indexing. Example lists: \url{https://www.ranks.nl/stopwords}.
  \item \textbf{When not to remove them}: some tools avoid removing stop words to
        support \emph{phrase search} (exact, not keyword-based): ``to be or not to be''
        consists almost only of stop words.
\end{itemize}

A word cloud of a novel (the slides show one for a Thomas Hardy novel) is dominated by
\emph{the}, \emph{and}, \emph{of}, \emph{to}, \emph{a}, \emph{was}, \emph{her}: they are
frequent in \emph{any} English text and say nothing about this one.

\subsubsection{Frequent in the document or frequent in the language?}
\slidesrc{Syntax: RE \& word: 24--29}

\begin{nbintuition}[the key question]
Is a word that appears often in a document always a stop word? No. A word can be
frequent in a document $d$ for two very different reasons:
\begin{itemize}[nosep]
  \item it is related to the \textbf{topic} of $d$ (``cat'' in a document about cats);
  \item it is frequent in the \textbf{language} (``the'' in every document).
\end{itemize}
To tell them apart, we must look not only at $d$ but at a whole \textbf{corpus}: a
topical word is frequent in \emph{some} documents, a stop word in \emph{all} of them.
\end{nbintuition}

The running example uses two short documents:

\begin{nbexample}{the corpus}
\textbf{Doc1} (54 words): \emph{The cat (Felis catus) is a small carnivorous mammal. It
is the only domesticated species in the family Felidae and often referred to as the
domestic cat to distinguish it from wild members of the family. The cat is either a
house cat or a farm cat, which are pets, or a feral cat.}

\textbf{Doc2} (42 words): \emph{A car (or automobile) is a wheeled motor vehicle used for
transportation. Cars became widely available in the early 20th century. One of the first
cars accessible to the masses was the 1908 Model T. Cars were rapidly adopted in the US.}

\begin{center}
\small
\begin{tabular}{@{}l c c c@{}}
\toprule
 & \textbf{``the''} & \textbf{``cat''} & \textbf{``car(s)''} \\
\midrule
Doc1 & 6 & 6 & 0 \\
Doc2 & 5 & 0 & 4 \\
\midrule
corpus $C_2 = \{\text{Doc1}, \text{Doc2}\}$ & 11 & 6 & 4 \\
\bottomrule
\end{tabular}
\end{center}
Corpus $C_1$ contains only Doc1; corpus $C_2$ contains Doc1 and Doc2.
\end{nbexample}

\begin{nbwarning}[a counting slip in the slides]
The slides count \textbf{4} occurrences of ``the'' in Doc2 (and 10 in $C_2$). Doc2
actually contains it \textbf{5} times: \emph{\underline{the} early}, \emph{\underline{the}
first}, \emph{\underline{the} masses}, \emph{\underline{the} 1908}, \emph{\underline{the}
US}. This changes $\text{tf}(\text{``the''}, \text{Doc2})$ from $4/42$ to $5/42$, but not
the final TF-IDF value, which is $0$ either way (see below).
\end{nbwarning}

In $C_1$ alone, ``the'' and ``cat'' are equally frequent (6 times each): with a single
document there is no way to tell a topical word from a stop word. In $C_2$ the
difference appears: ``the'' is in \emph{both} documents, ``cat'' only in one.

\subsubsection{Term frequency and inverse document frequency}
\slidesrc{Syntax: RE \& word: 30--34}

\begin{nbdef}{TF-IDF}
\textbf{Term frequency}: the number of times term $t$ occurs in document $d$, usually
normalised by the length of the document:
\[
  \text{tf}(t, d) = \frac{\text{number of occurrences of } t \text{ in } d}{\text{length}(d)}.
\]
\textbf{Inverse document frequency}: the logarithm of the inverse fraction of documents
of the corpus $C$ that contain $t$:
\[
  \text{idf}(t, C) = \log \frac{\abs{C}}{\abs{\{d \in C : t \in d\}}},
\]
where $\abs{C}$ is the number of documents in the corpus.

\textbf{TF-IDF}: the product of the two,
\[
  \text{tf-idf}(t, d, C) = \text{tf}(t, d) \cdot \text{idf}(t, C).
\]
\end{nbdef}

\begin{itemize}
  \item \textbf{tf} is high when $t$ is frequent \emph{in this document}.
  \item \textbf{idf} is high when $t$ is \emph{rare in the corpus}. If $t$ occurs in
        every document, the fraction is $\abs{C}/\abs{C} = 1$ and $\text{idf} = \log 1 = 0$:
        the word cannot distinguish documents at all, exactly like a stop word.
  \item The \textbf{logarithm} dampens the effect: a word in 1 document out of 1000 is
        more informative than one in 10 out of 1000, but not 10 times more.
  \item If $t$ appears in \emph{no} document the denominator is 0. A common fix is to
        add 1 to the denominator; it makes sense when the corpus is large enough, since
        it slightly distorts the values.
\end{itemize}

\begin{nbwarning}[base of the logarithm]
The slides' numbers use the base-10 logarithm ($\log 2 \approx 0.301$). Other sources use
the natural logarithm ($\ln 2 \approx 0.693$). The base only rescales all idf values by
the same factor, so the \emph{ranking} of words does not change; just be consistent.
\end{nbwarning}

\subsubsection{Worked example on \texorpdfstring{$C_2$}{C2}}
\slidesrc{Syntax: RE \& word: 35--46}

\begin{nbexample}{TF-IDF of ``the'', ``cat'', ``car''}
First the idf values ($\abs{C_2} = 2$):
\begin{align*}
  \text{idf}(\text{the}, C_2) &= \log\tfrac{2}{2} = \log 1 = 0
      \reason{``the'' is in both documents}\\
  \text{idf}(\text{cat}, C_2) &= \log\tfrac{2}{1} = \log 2 \approx 0.301
      \reason{only in Doc1}\\
  \text{idf}(\text{car}, C_2) &= \log\tfrac{2}{1} = \log 2 \approx 0.301
      \reason{only in Doc2}
\end{align*}
Then the products:
\begin{align*}
  \text{tf-idf}(\text{the}, \text{Doc1}, C_2) &= \tfrac{6}{54}\cdot 0 = 0 \\
  \text{tf-idf}(\text{the}, \text{Doc2}, C_2) &= \tfrac{5}{42}\cdot 0 = 0 \\
  \text{tf-idf}(\text{cat}, \text{Doc1}, C_2) &= \tfrac{6}{54}\cdot 0.301 \approx 0.111 \cdot 0.301 \approx 0.033 \\
  \text{tf-idf}(\text{cat}, \text{Doc2}, C_2) &= 0 \cdot 0.301 = 0
      \reason{tf $=0$, idf irrelevant}\\
  \text{tf-idf}(\text{car}, \text{Doc1}, C_2) &= 0 \cdot 0.301 = 0 \\
  \text{tf-idf}(\text{car}, \text{Doc2}, C_2) &= \tfrac{4}{42}\cdot 0.301 \approx 0.095 \cdot 0.301 \approx 0.029
\end{align*}
\end{nbexample}

\begin{center}
\small
\begin{tabular}{@{}l c c c@{}}
\toprule
\textbf{tf-idf} & \textbf{``the''} & \textbf{``cat''} & \textbf{``car''} \\
\midrule
Doc1 & 0 & \textbf{0.033} & 0 \\
Doc2 & 0 & 0 & \textbf{0.029} \\
\bottomrule
\end{tabular}
\captionof{table}{TF-IDF in $C_2$. Each document's only non-zero word is its topic.}
\end{center}

``The'' is as frequent as ``cat'' in Doc1, but its TF-IDF is $0$ because it appears
everywhere; ``cat'' and ``car'' come out as the words that characterise their documents.
This is how TF-IDF finds \emph{keywords} and, symmetrically, \emph{stop words} (terms with
idf close to 0) automatically. Online analysers such as
\url{https://www.online-utility.org/text/analyzer.jsp} compute word frequencies of a
text you paste.

\begin{nbwhy}
TF-IDF is one of the text features of Section~\ref{sec:course-map}: representing each
document by the TF-IDF values of its words gives a vector that search engines use to
rank documents for a query and classifiers use to decide the topic of a text.
\end{nbwhy}
